%% Resumo atualizado após comparação com a implementação (12/11/2026)
\begin{resumo}

O aumento de pessoas que vivem sozinhas motiva a criação de ferramentas de apoio à verificação periódica de bem-estar. Este trabalho descreve o desenvolvimento do Linha Vital, protótipo acadêmico de aplicativo Android integrado a uma API Kotlin/Spring Boot e banco de dados PostgreSQL. O sistema permite cadastro e autenticação de usuários, gerenciamento de contatos priorizados e configuração de intervalos de monitoramento. O funcionamento preventivo baseia-se em confirmações explícitas de bem-estar (\textit{check-ins}): após o vencimento do intervalo e de uma tolerância configurada no servidor, o backend registra a ocorrência e tenta encaminhar alertas por e-mail aos contatos cadastrados. Em situação de emergência percebida pelo próprio usuário, o aplicativo oferece acionamento manual de SOS e sequência assistida de ligações, com registro do resultado informado pelo usuário. A localização é obtida pontualmente, mediante permissão, e pode complementar a comunicação quando disponível. O projeto compreendeu levantamento de requisitos, prototipação, implementação e validação funcional. Em 14 de outubro de 2026, uma bateria pré-deploy documentou 19 cenários funcionais, dos quais 17 foram aprovados e dois aprovados com observações, com 111 capturas de evidência; os ensaios registraram o recebimento de e-mail no ambiente de teste, sem representar garantia de entrega em produção. Também foram identificados testes automatizados no repositório, mas não foram aferidas cobertura de código, disponibilidade ou desempenho sob carga. O resultado é um MVP demonstrável para fins acadêmicos, cuja implantação pública depende de avaliação adicional de segurança, usabilidade, confiabilidade e conformidade legal. O Linha Vital não substitui serviços profissionais de emergência.

\vspace{\onelineskip}
\noindent
\textbf{Palavras-chave}: check-in; monitoramento de inatividade; aplicativo Android; contatos de emergência; sistema de apoio.
\end{resumo}
